% Folder 34, batch 02 — archivists' pages 21–40 (author's pp. 20–39).
% Pass: Opus 5.5 (claude-opus-5-5), 2026-10-03 — first pass, unchecked against the pages by a human.
% Content: quasi-unipotent l-adic representations (§5, end), residue field of finite
% type (§6), the monodromy theorem as conjecture 7.1 (§7) and its arithmetic proof (§8).
\documentclass[11pt,a4paper]{article}
\input{../preamble/grothendieck}

\folder{34}
\batch{2}
\pages{21}{40}
\foldertitle{SGA 7 : notes manuscrites (s.d.).}
\dating{[vers 1967-1973]}
\watermark{Édition de démonstration}

\begin{document}

\page{21}\note{p. 20 de l'auteur ; la phrase commence à la page précédente, hors du lot.}
la représentation \emph{quasi-unipotente} s'il existe un sous-groupe ouvert $U$ de $\pi$ tel que la représentation induite $u|U$ soit unipotente.

5.2. Supposons $R$ infini, et soit $G$ le groupe \uncertain{algébrique} engendré par $u(\pi)$, \struck{Alors} égal à l'adhérence (pour la structure induite réduite) de l'ens.\ $u(\pi)$\note{la page porte un signe que l'on peut lire $u(G)$ ; le sens demande $u(\pi)$.} dans le schéma en groupes des \struck{\ill{}} \add{automorphismes} de $E$. Comme l'ens.\ des \struck{\ill{}} \add{automorphismes} unipotents \add{est} une partie Zariski-fermée de l'ens.\ $\mathrm{Aut}(E)$, il s'ensuit que $u$ est unipotente ssi $G$ est unipotent (i.e.\ \uncertain{tous} ses éléments sur $\overline{R}$ sont des automorphismes unipotents de $E_{\overline{R}}$).

\marginal{Notons que $G/G^{0}$ est un groupe \uncertain{discret}, isom.\ à $\pi/U$, si $U$ est l'image inverse de $G^{0}(R)$.}

5.3. \struck{Dès lors}, Si $R$ est un corps topologique, et $u$ est continue, alors dire que $u$ est quasi-unipotente revient à dire que $G^{0}$ est unipotent. C'est évidemment suffisant, \add{par la top.\ discrète de $G/G^{0}$}, car l'homomorphisme
\[ \pi \longrightarrow (G/G^{0})(R) \simeq G(R)/G^{0}(R) \]
est continu, donc le noyau \add{l'image inverse} de $G^{0}(R)$ est un sous-groupe ouvert $U$ de $\pi$. Inversement, si $u$ est quasi-unipotente, il y a un sous-groupe ouvert $U$ de $\pi$ qui opère de façon unipotente, prenons $U$ assez petit \ill{} $U \subset u^{-1}(G^{0}(R))$. Alors $\overline{u(U)}$ est un sous-groupe algébrique \add{(lisse $H$)} de $G$, dont

\marginal{\drawing{petit croquis à l'encre et au crayon dans la marge inférieure gauche : des arcs, des traits, la lettre $R$ et le mot \uncertain{constante} souligné ; le croquis ne se laisse pas lire.}}

\page{22}\note{p. 21 de l'auteur.}
la contient immédiatement qu'il \uncertain{suffit} \struck{\ill{}} \add{il} \ill{} $H(\overline{R})$ d'indice fini dans $G^{0}(\overline{R})$, \uncertain{implique} que $H = G^{0}$ ; \struck{donc} comme $H$ est unipotent, il en est de même de $G^{0}$.

\marginal{dès qu'on \uncertain{dispose de} la condition $E^{\struck{H}} = E^{G^{0}}$.}

5.4. Supposons maintenant que $\pi$ opère de façon quasi-unipotente (i.e.\ $G^{0}$ est unipotent) et supposons que $R$ est de car.\ nulle. \struck{Si $U$ est} \ill{} \struck{\ill{}}

\note{Le passage qui suit, jusqu'à la fin de la page, est encadré et barré de deux longs traits obliques ; il est transcrit tel qu'il se lit.}
\struck{Soit $U$ \ill{} $\pi$ qui opère de façon unipotente (\ill{} $u^{-1}(G^{0}(R))$), \ill{} $E^{U}$, sur lequel $\pi/U$ opère. Je dis que $E^{U}$ ne dépend pas du choix de $U$ : il suffit de voir que $E^{U} = E^{U'}$ si $U' \subset U$, $U'$ invariant. Or $\pi/U'$ opère sur $E^{U'}$, donc via un groupe fini qui opère de façon unipotente, donc trivialement. De plus, $E^{U}$ \ill{} ne dépend pas du sous-groupe $U$ invariant de $\pi$). Ceci nous permet de définir une filtration canonique de $E$, en posant}
\[ \text{\struck{$E^{(0)} = 0$, \quad $E^{(i+1)} = $ image inverse de $(E/E^{i})^{U} = (E/E^{i})^{G^{0}}$ dans $E$.}} \]
\struck{L'hypothèse que $G^{0}$ opère de façon unipotente implique que cette filtration est exhaustive. D'autre part, $U$ opère \uncertain{trivialement} sur les gradués associés, par construction.}

\page{23}\note{p. 22 de l'auteur.}
Alors je dis que \emph{si} $g \in G(R)$, $g$ est \emph{unipotent} ssi $g \in G^{0}(R)$. En effet, \struck{on est réduit à prouver seulement \og seulement si \fg}. \add{Il suffit de prouver} : $g$ unipotent $\Rightarrow g \in G^{0}(R)$. On \struck{considère une} l'image de $g$ dans le groupe quotient $\Gamma = G/G^{0}$ est un \uncertain{élément} unipotent de $\Gamma(R)$, et comme \ill{} est un \ill{}, c'est l'unité. \struck{Par \ill{} Autrement dit} :

\struck{On filtre $E$ par}
\note{Le passage qui suit est encadré, et barré d'un trait oblique.}
\struck{$E^{0} = 0$, … $E^{i}$ …, $\xi$ avec}
\[ \text{(5.4.1)} \qquad E^{i+1} = \text{image inverse dans $E$ de } (E/E^{i})^{G^{0}} . \]
Cette filtration est \add{exhaustive, i.e.\ $\exists\, i$ t.q.\ $E^{i} = E$}, \struck{finie} et \uncertain{on a} \uncertain{explicitement} \ill{} \ill{} simplement que $G^{0}$ est unipotent. Alors $G/G^{0}$ opère sur les $E^{i+1}/E^{i}$, et si $g \in G$ est unipotent, il induit des op.\ unipotents sur les $E^{i+1}/E^{i}$, donc des opérations triviales, i.e.\ il \uncertain{provient} d'un \uncertain{élément} du \ill{} $\mathrm{Gr}(E)$. Mais comme \ill{} \ill{} \uncertain{de} l'op.\ de $G/G^{0}$ est \uncertain{fidèle} \struck{\ill{}} \ill{}, \struck{$E \sim \mathrm{Gr}(E)$ par une} \ill{}.

En particulier, l'image inverse $U$ de $G^{0}(R)$ dans $R$ \uncertain{est} le plus grand sous-groupe de $\pi$ opérant de façon unipotente, \uncertain{a fortiori} le plus grand sous-groupe ouvert ayant cette propriété.

5.5.\ \struck{\ill{}} Notons que \struck{\ill{}} les \uncertain{composantes} connexes des groupes \uncertain{algébriques} \add{unipotents} commutatifs \add{sur un corps $R$ de car.\ nulle} sont connexes, \uncertain{prenant} les filtrations \og caractéristiques \fg\ de $G^{0}$ (par les commutateurs successifs, $G^{(0)} = G$,

\page{24}\note{p. 23 de l'auteur.}
$G^{(i+1)} = [G^{(i)}, G]$) tels que $\mathrm{Gr}(G^{0})$ soit un groupe vectoriel. On en conclut, comme
\[ H^{i}(\Gamma, V) = 0 \quad \text{pour } i \geq 0 \]
si $V$ est vectoriel sur lequel $\Gamma$ opère (car $V$ est uniquement divisible …), que l'\emph{extension} $G$ de $\Gamma$ par $G^{0}$ \emph{est \uncertain{splittée} de façon unique à un unique} \ill{} près, \uncertain{ce qui} \uncertain{veut dire}
\[ \text{(5.5.1)} \qquad G \simeq \Gamma \cdot G^{0} , \]
\uncertain{Ceci} \struck{\ill{}} \add{$\Gamma$ opérant de façon naturelle} sur $E$, \struck{et \ill{} \ill{} \ill{} $\Gamma$ \ill{} \ill{} que $g \in G(R)$, \ill{}} \struck{$g$ unipotent $\Rightarrow g \in G^{0}(R)$}].
\ill{}

\struck{(5.6)} : \struck{5.6} \quad P.\ ex.

\page{25}\note{p. 24 de l'auteur ; la page est barrée de deux longs traits obliques.}
5.5. On en conclut que \struck{il existe le plus grand} \add{l'ens.\ des} $g \in \pi$ tels que $u(g)$ \struck{soit} unipotent est un sous-groupe ouvert invariant de $\pi$, à savoir $u^{-1}(G^{0}(R))$ ; c'est évidemment le plus grand sous-groupe \add{(ouvert)} \struck{\ill{}} de $\pi$ qui opère de façon unipotente. \add{Notons-le $U$}. On a
\[ \text{(5.5.1)} \qquad \pi/U \simeq (G/G^{0})(R) = G(R)/G^{0}(R) . \]
Soit $U'$ un sous-groupe \add{ouvert} \struck{\ill{}} de $\pi$, \struck{\ill{}} opérant de façon unipotente, i.e.\ un sous-groupe \uncertain{ouvert} de $U$, je dis

\[ E^{U'} = E^{U} = E^{G^{0}} . \]
En effet, on voit \uncertain{supposer} que $U'$ est invariant, \add{alors} $U/U'$ opère sur $E^{U'}$, et \struck{\ill{}} \add{c'est} un groupe fini qui opère de façon unipotente, donc trivialement, donc $E^{U'} = E^{U}$. On a $E^{U} = E^{G^{0}}$, car il suffit de \ill{} que $\overline{u(U)} = G^{0}$.

5.6. Supposons maintenant que $R = \mathbf{Q}_{\ell}$, \struck{alors} \add{(et $G^{0}$ commutatif)}, \add{i.e.\ que l'opération de $\pi$ sur $E$ soit \og quasi-déployée \fg}. \add{Alors}
\[ U \longrightarrow G^{0}(\mathbf{Q}_{\ell}) \simeq \mathbf{Q}_{\ell}^{d} , \]
\marginal{Donc $G^{0} \simeq \mathbf{G}_{a}^{d}$ (avec \uncertain{structure vectorielle canonique})}
comme l'image un sous-groupe compact de $\mathbf{Q}_{\ell}^{d}$, lequel est \uncertain{nécessairement} isomorphe à un $\mathbf{Z}_{\ell}^{d'}$ ($d' \leq d$). Comme \struck{\ill{}} $\overline{u(U)}$ engendre le groupe algébrique $G^{0}$, il s'ensuit aussitôt que $d' = d$, \uncertain{et} que $u(U)$ est un sous-$\mathbf{Z}_{\ell}$-module de $G^{0}(\mathbf{Q}_{\ell})$ qui engendre l'esp.\ vectoriel $G^{0}(\mathbf{Q}_{\ell})$.

\page{26}\note{p. 25 de l'auteur.}
Si
\[ \text{(5.6.1)} \qquad N = \operatorname{Ker}\bigl(u : \pi \longrightarrow G(\mathbf{Q}_{\ell}) \subset \mathrm{Aut}(E)\bigr), \quad \text{(donc $N \subset U$)}, \]
on voit donc que
\[ U/U \cap N \hookrightarrow G^{0}(R) \]
s'identifie à un $\mathbf{Z}_{\ell}$-module libre de rang $d$ \add{dans} $G^{0}(R)$ \add{donc} \uncertain{engendrant} \uncertain{rationnellement} $G^{0}(R)$ qui engendre $\mathbf{Q}_{\ell}$. D'ailleurs, on trouve ainsi un \uncertain{homomorphisme} d'extensions

\begin{tikzcd}[column sep=small]
1 \arrow[r] & U/N \arrow[r] \arrow[d, hook] & \pi/N \arrow[r] \arrow[d, hook] & \pi/U \arrow[r] \arrow[d, "\wr"] & 1 \\
1 \arrow[r] & G^{0}(R) \arrow[r] & G(R) \arrow[r] & G(R)/G^{0}(R) \arrow[r] & 1
\end{tikzcd}

Ainsi, le \struck{\ill{}} groupe algébrique $G(R)$ est \struck{\ill{}} canoniquement déterminé \add{à iso unique près}\note{la parenthèse ajoutée, sous la ligne, se lit aussi \uncertain{\og bornée \fg} ; la lecture reste douteuse.} (à \uncertain{isomorphisme} \add{unique} près) par la connaissance du sous-groupe $U$ de $\pi$, et du \add{fermé} sous-groupe $N$ de $\pi$ tel que $U/N \simeq \mathbf{Z}_{\ell}^{d}$. $G(R)$ est en effet l'image de l'extension canonique de $\pi/U$ par $U/N$ \struck{\ill{}} par l'homomorphisme canonique
\[ U/N \longrightarrow U/N \otimes_{\mathbf{Z}_{\ell}} \mathbf{Q}_{\ell} . \]

[N.B.\ : la structure de $\mathbf{Z}_{\ell}$-module sur le groupe \uncertain{profini} $U/N$ est canoniquement déterminée), et la structure de groupe algébrique se détermine aussi trivialement …

Si ${}^{1}M_{\ell} = U/[U,U]^{(\ell)}$ \add{\uncertain{complété} $\ell$-primaire} est un $\mathbf{Z}_{\ell}$-module de t.f.\ \add{l'extension correspondante}, alors on construit \uncertain{une} \struck{\ill{}} \add{\ill{}} \uncertain{extension} universelle de $\Gamma = \pi/U$ par le groupe vectoriel \add{\uncertain{correspondant}} : $M_{\ell} \otimes_{\mathbf{Z}_{\ell}} \mathbf{Q}_{\ell}$.

\page{27}\note{p. 26 de l'auteur.}
5.7. Comme groupe algébrique abstrait cependant, $G$ est déterminé par la connaissance de $\Gamma \simeq \pi/U$, et de sa représentation par automorphismes du vectoriel \struck{$V$} $E$ \struck{\ill{}} du vectoriel
\[ U/N \otimes_{\mathbf{Z}_{\ell}} \mathbf{Q}_{\ell} , \]
quotient du vectoriel
\[ \underbrace{U/[U,U]^{(\ell)}}_{\text{complété $\ell$-primaire}} \otimes_{\mathbf{Z}_{\ell}} \mathbf{Q}_{\ell} . \]

Lorsque $\pi = I \simeq \mathbf{Z}_{\ell}(1)$ \add{$\mathrm{Gal}(\overline{K}/K) \simeq$} comme on sait, \add{(cas \uncertain{strict.} local)}, $\Gamma$ puisque $U$ est d'indice fini, $\Gamma_{U}$\note{lecture douteuse : un $\Gamma$ indicé d'un $U$ souligné deux fois, à côté de « iso ».} iso\ill{}, \uncertain{\ill{}} isomorphisme canonique
\[ \text{(5.7.1)} \qquad U/[U,U]^{(\ell)} \simeq \mathbf{Z}_{\ell}(1) \]
(isom.\ canonique), \struck{donc \ill{} \ill{}} \add{et} \add{une} représentation naturelle de $\pi/U = \Gamma$ « dans $\mathbf{Z}_{\ell}(1)$ », i.e.
\[ \pi/U = \Gamma \longrightarrow \mathbf{Z}_{\ell}^{*} \]
qui est connue pour être \emph{triviale} \add{(3.15)}. Donc on aura un isom.\ \add{canonique} déduit de (5.5.3)\note{lu « (5.5.3) » ; aucune formule de ce numéro n'apparaît dans le lot.}
\[ \text{(5.7.2)} \qquad G \simeq G^{0} \times \Gamma \quad \text{\struck{\ill{}}} \]
\struck{et $G^{0}$ est le groupe \ill{}} \emph{i.e.\ ici $G^{0}$ est soit nul, soit} \struck{\ill{}} \emph{canoniquement isomorphe au} \struck{$\mathbf{Q}_{\ell}(1)$}, \emph{groupe algébrique défini par le vectoriel} $\mathbf{Q}_{\ell}(1)$. On a en particulier
\[ \text{(5.7.3)} \qquad \begin{cases} G(\mathbf{Q}_{\ell}) \simeq \Gamma = \pi/U & \text{si $G$ fini, i.e.\ $U$ opère trivialement,} \\ G(\mathbf{Q}_{\ell}) \simeq \mathbf{Q}_{\ell}(1) \times \Gamma & \text{sinon.} \end{cases} \]

\page{28}\note{p. 27 de l'auteur. Le haut de la page (5.10 et 5.10.1) est barré de grands traits croisés ; la section 6 qui suit, de deux traits obliques.}
\struck{5.10. On peut aussi regarder, \ill{} \ill{} pour une représentation $E$ \uncertain{$\ell$-adique} donnée de $I$, le sous-module}
\[ \text{\struck{(5.10.1) \quad $E^{P}$}} \]
\struck{des invariants sous $P$. Il est stable par $I$, et par passage au quotient, $I/P = I_{t} \simeq \prod_{\ell \neq p} \mathbf{Z}_{\ell}(1)$ opère sur $E^{P}$.}

\struck{En tout cas, une représentation $\ell$-adique quasi-unipotente de $\pi$ \uncertain{s'exprime} \uncertain{en termes} \uncertain{d'une} représentation \emph{algébrique} de $\Gamma \times_{\mathbf{Q}_{\ell}} W(\mathbf{Q}_{\ell}(1))$}.
\marginal{On trouve des systèmes projectifs de telles \ill{} sur $\mathbf{Q}_{\ell}$ : $\Gamma_{i} \times W(\mathbf{Q}_{\ell}(1))$ ($\Gamma_{i}$ quotients $I/V_{i}$, $V_{i}$ sous-groupes \uncertain{ouverts} \uncertain{invariants}), \ill{} les « \uncertain{variétés} » donnant une \emph{$\otimes$-catégorie équivalente à celle des repr.\ $\ell$-adiques de $I$ quasi-unipotentes}.}

\section{6. Cas d'un corps résiduel fini}
\note{titre de l'auteur, souligné, barré avec le reste de la page ; « Cas \struck{\ill{}} d'un corps \uncertain{résiduel} fini ».}

\marginal{\uncertain{généralisation} à un corps résiduel de type fini \ill{} \ill{} \ill{} obtenir \ill{} \ill{} \ill{} l'exigence \ill{}\note{note de marge écrite en oblique, à moitié couverte par les traits ; lecture très incertaine.}}

\struck{Si $k$ …}
\[ \text{\struck{(6.1) \quad card $k = q = p^{f}$}} \]
\struck{on sait que}
\[ \text{\struck{$\pi_{0} \simeq \widehat{\mathbf{Z}}$ ,}} \]
\struck{(6.2) l'isomorphisme étant défini par le $\mathrm{frob}_{q}$ de $\pi_{0}$, défini par}
\[ \text{\struck{(6.3) \quad $\mathrm{frob}_{q}(\lambda) = \lambda^{q}$ \quad ($\lambda \in \overline{k}$) .}} \]
\struck{Cet automorphisme opère sur $\overline{k}^{*}$ par élévation à la puissance $q$, donc il opère de même sur les $\mu_{n}(\overline{k}^{*})$, donc sur $I_{t} = \prod_{\ell \neq p} \mathbf{Z}_{\ell}(1)$.}

\struck{Par suite, \ill{} \ill{} se donner une représentation \ill{} \ill{} (\ill{} quelconque $R$) donnée \ill{}}
\[ \text{\struck{$u : \pi \longrightarrow \mathrm{Aut}_{R}(E)$}} \]
\struck{est \ill{} triviale sur le \ill{} $I^{P}$…, \ill{} donc, considérons $I_{t}$ comme \ill{}, \ill{} isomorphisme : \ill{} quotient de}
\[ \text{\struck{$\widehat{\mathbf{Z}} \simeq \prod_{\ell} \mathbf{Z}_{\ell}$}} \]

\page{29}\note{p. 28 de l'auteur. La section 6 reprend ici, sous un titre corrigé, ce que la p. 28 du fonds avait commencé et barré.}
\section{6. Cas d'un corps résiduel de type fini}

\textbf{Proposition 6.1.} \emph{Soit $\pi$ un groupe profini, extension d'un groupe profini $\pi_{0}$ par un groupe profini} \struck{\ill{}} \add{\emph{résoluble}} \struck{\ill{}} \emph{$I$, \add{(\uncertain{supposons} que} pour tout sous-groupe ouvert $\pi'_{0}$ de $\pi_{0}$, \uncertain{…}) \add{(\ill{} \ill{} $\pi'$ \ill{} sous-groupe ouvert de $\pi$, $I' = \pi' \cap I$,} [N.B.\ si $I$ \uncertain{est} \struck{\ill{}} \add{commutatif}, \add{\uncertain{il} \uncertain{suffit}} $I'_{\pi'_{0}} = \varprojlim (I'_{\alpha})_{\pi'_{0}}$ \struck{$= I'/[I',\pi']$} … (i.e.\ $I$ prenant les quotients \struck{finis} de $I$ par un sous-groupe ouvert, invariant dans $\pi$)] \struck{et $\pi'$ un groupe} \struck{invariant dans $\pi$} [fini] \struck{i.e.} \add{posant $I' = \pi' \cap I$,}}
\[ [I', \pi'] \quad \text{\emph{est d'indice fini dans} } I' , \quad \text{\emph{pour tout $\pi'$ ouvert}} \]
\emph{\uncertain{invariant} de $\pi'_{0}$ dans $\pi$}). \emph{Alors pour toute représentation $\ell$-adique de $\pi$, la représentation induite sur $I$ est quasi-unipotente.}

\marginal{\uncertain{Vérifier} \uncertain{nouvelle} \uncertain{version} à cause des \uncertain{groupes} …\note{note de marge écrite en oblique, à la hauteur de l'énoncé ; lecture très incertaine.}}

Soit $E$ le vectoriel de dimension finie sur $\mathbf{Q}_{\ell}$ sur lequel on fait opérer $\pi$. Soit \struck{$G = \ldots ($} $G(H)$ la composante neutre du groupe algébrique \add{\uncertain{engendré} par} $\pi$ (resp.\ $I$), on sait que lorsque \struck{\ill{}} \add{$\pi$ \uncertain{résoluble}} défini par $\pi$ (resp.\ $I$) dans $G$. L'hypothèse \struck{\ill{} signifie} \add{sur $\pi$ implique})
\[ H \subset G , \quad H \text{ invariant dans } G , \]
que
\[ [G, H] = H . \]
On a aussi : $H = \struck{H_{\ill{}}\, H_{u}}$ \add{$T \cdot U$} (\uncertain{produit semi-direct} de \add{partie} \struck{\ill{}} \add{torique} et partie unipotente), \struck{$H_{r}$ et $H_{u}$ \ill{}}, \add{$U$, donc $H/U \simeq T$} et $G$ opère sur \struck{$H_{r}$} $T$ \struck{(\ill{} \ill{} \ill{})} trivialement (\uncertain{$G$ étant} \ill{} connexe). \uncertain{Alors} $[G, T] = T$ donc $T = 1$,
\[ [G, H] = H \quad \text{implique} \quad \text{\struck{$H_{r} = 1$,}} \]
donc $H$ est unipotent, cqfd.

\page{30}\note{p. 29 de l'auteur.}
\textbf{Proposition 6.2.} \emph{Soit \struck{$\pi = \mathrm{Gal}(\overline{K}/K)$}, $\pi$ comme dans 6.1, \struck{\ill{} le corps $k$} supposons que le corps résiduel $k$ soit tel que l'extension \add{\uncertain{obtenue}} de $\overline{k}$ engendrée par les racines $\ell^{\nu}$-ièmes de l'unité ($\nu \in \mathbf{N}$) soit infinie. Alors \struck{toute extension} \uncertain{toute} représentation $\ell$-adique de $\pi$ induit sur $I$ une repr.\ quasi-unipotente.}

En effet, quitte à remplacer $K$ par une extension finie $K'$ de $K$, et $I$ par un sous-groupe ouvert, on peut supposer que l'action de $I$ se fait à travers le quotient $I(\ell) = I_{t}(\ell) = \mathbf{Z}_{\ell}(1)$, \uncertain{\ill{}} \uncertain{l'hyp.}\ de 6.1 se \uncertain{traduit} par \ill{} … : \uncertain{la} \uncertain{représentation} d'un quotient $\pi'$ de $\pi$, \ill{} \uncertain{extension} de $\pi_{0}$ par $\mathbf{Z}_{\ell}(1) \simeq \mathbf{Z}_{\ell}$. L'hypothèse faite sur $k$ signifie que l'hyp.\ de 6.1 est vérifiée.

\textbf{Corollaire 6.3.} \emph{Toute représentation $\ell$-adique \struck{\ill{}} \add{\uncertain{continue}} de $\pi$ induit une repr.\ \uncertain{unipotente} \uncertain{triviale} de $I$.}

\uncertain{Cf.}\ II 1.12, qu'il y aurait lieu de \uncertain{reprendre} ici.

\page{31}\note{p. 30 de l'auteur. En tête, « I » puis « 6. » : une nouvelle rédaction de la proposition 6.1 commence.}
\textbf{Prop.\ 6.1.} \emph{Soit}
\[ 1 \longrightarrow I \longrightarrow \pi \longrightarrow \pi_{0} \longrightarrow 1 \]
\emph{une extension de groupes profinis, et supposons que pour tout sous-groupe ouvert $\pi'$ de $\pi$, posant $I' = I \cap \pi'$, le} \struck{\emph{l'action de $I'$ sur \ill{} \ill{} \ill{} $I'$, i.e.\ le groupe $[I', \pi']$ soit ouvert dans}} \emph{\add{$I'^{\mathrm{ab}}(\ell)_{\pi_{0}}$} plus grand quotient \uncertain{pro-fini} de $I'^{\mathrm{ab}}(\ell)$, sur lequel $\pi'_{0} = \pi'/I'$ opère trivialement, est un groupe fini. Soit}
\[ u : \pi \longrightarrow \mathrm{Aut}_{\mathbf{Q}_{\ell}}(E_{\ell}) \]
\emph{une représentation $\ell$-adique de $\pi$,} \struck{\emph{soit et}} \emph{considérons} \struck{\emph{la représentation induite par $u$ sur}} \emph{$H$ ($G$) le sous-groupe algébrique de $\mathrm{Aut}_{\mathbf{Q}_{\ell}}(E_{\ell})$ engendré par $u(I)$ (resp.\ $I(G)$). Alors}

\emph{a) $H$ est invariant dans $G$, donc $H^{0}$ invariant dans $G$} \struck{\emph{et \ill{}, dans $G^{0}$}}.

\emph{b)} \struck{\emph{L'affirmation de $G^{0}$ sur \ill{}}} \emph{Soit $R$ le radical de $H^{0}$,} \add{\emph{(quotient inv.\ dans $G$, et sur}} \emph{$A = R/[R,R]$,} \emph{\add{Alors} l'action de $G^{0}$ sur $A$} \struck{\emph{satisfait \ill{} \ill{} \ill{}}} \struck{\emph{sont \ill{}}} \emph{\ill{} \ill{} la condition}
\[ A_{(G^{0})} = 0 , \quad \text{\emph{i.e.\ $A$ est} \struck{\emph{engendré}} \emph{engendré par les} } (1-g)(A), \ \text{\emph{où} } g \in G^{0}(\mathbf{Q}_{\ell}) . \]

\emph{c) $R$ est un groupe \uncertain{unipotent}.}

\textbf{Dém.} a) Résulte trivialement du fait que $I$ invar.\ dans $\pi$. Je suppose que $R$ est aussi invar.\ dans $G$, donc $G$ opère sur $A = R/[R,R]$.

b) Exprime l'hypothèse $I'_{\mathrm{ab}}(\ell)$ fini, en notant que $H^{0}/[H^{0},H^{0}] \longleftarrow R/[R,R] = A$ est une isogénie …

c) A priori, $R = T \cdot U$, $U$ unipotent, $T$ tore ($U$ ss-groupe car.\ de $R$). Le groupe $U$ est invariant \struck{dans $R$, et $G$ opère} sur $R/U = T$, de façon que $G^{0}$ opère trivialement (car le groupe des aut.\ d'un tore \uncertain{fermement} …

\page{32}\note{p. 31 de l'auteur ; le bas de la page est vide.}
un groupe discret). Donc $T$ est un quotient de $A_{G^{0}}$, donc est nul ; i.e.\ $R = U$ est unipotent.

\textbf{Corollaire 6.2.} \emph{Supposons que le \uncertain{radical} … résoluble. Alors} \struck{\ill{}} \emph{$u|I$ est \uncertain{ess.}\ unipotent.}\note{la ligne précédente se lit « que la restriction $u|I$ soit ess.\ v. » ; l'énoncé est très abrégé.}

\textbf{Prop.\ 6.3.} Application aux groupes de monodromie …

\page{33}\note{p. 32 de l'auteur.}
\struck{Ce cas, on conclut grâce à 6.6.}

\section{7. Le théorème de monodromie}
\note{titre de l'auteur, souligné ; à sa suite, une ligne biffée où se lit \emph{« énoncé conjectural »} au-dessus d'un mot raturé.}

\struck{L'énoncé est le suivant} \add{(Serre-Tate)}.

\textbf{Conjecture 7.1.} \emph{\struck{Soit $X$ un schéma de type fini sur $K$,} \add{(Soit $X$ …)} \struck{\ill{} \ill{}, … $\neq p$, et \ill{}} \struck{\ill{} supposons que $H^{i}(X_{\overline{K}}, \mathbf{Z}/\ell\mathbf{Z})$ est fini} \add{cf.\ 0.3)}. Sous les conditions de l'exemple 0.3, et supposons $\ell \neq p$, la représentation induite sur le groupe d'inertie}
\[ I \longrightarrow \mathrm{Aut}_{\mathbf{Z}_{\ell}}\bigl(H^{i}(X_{\overline{K}}, \mathbf{Z}_{\ell})\bigr) \]
\emph{est quasi-unipotente, i.e.\ $\exists$ sous-groupe d'indice fini $U$ dans $I$ tel que l'image de $U$ soit formée d'opérations unipotentes dans $H^{i}(X_{\overline{K}}, \mathbf{Z}_{\ell})$.}

7.2. Nous allons voir par la suite deux démonstrations \uncertain{\og relatives \fg} de cette conjecture, utilisant l'une et l'autre \struck{\ill{} \ill{} \ill{}} la résolution des \emph{singularités}, de façon d'ailleurs assez différente. De ces démonstrations, il résulte que la conjecture est vraie dans chacun des cas suivants :
a) $p = 1$, i.e.\ $k$ (ou \uncertain{ce} \uncertain{qui} \uncertain{revient} au même, $K$) de car.\ nulle ;
b) $i = 1$, \marginal{\struck{\ill{}}\note{mot de la marge biffé et encerclé, relié à b) par un trait.}}
c) $S$ est le hensélisé (ou le \struck{\ill{}} \add{hensélisé strict}) d'un anneau local d'un point d'un schéma de type fini sur un corps, ou sur $\mathrm{Spec}\,\mathbf{Z}$, et $X_{K}$ est propre sur $K$ \emph{ou} $K$ de car.\ nulle.
La première démonstration, qui est de nature arithmétique, sera donnée dans la

\page{34}\note{p. 33 de l'auteur.}
présent exposé en se réduisant via 5.6 \struck{\ill{}} au cas d'un corps résiduel fini. Cette \uncertain{démonstration} \uncertain{avec} \uncertain{une} \uncertain{résolution} \struck{s'applique} \add{pour donner la conj.\ dans les cas a) et c)}. \struck{Dans} La démonstration géométrique fera l'objet de l'exposé suivant ; elle permettra de préciser assez considérablement \struck{\ill{}} la conclusion de unipotence, et \struck{\ill{}} s'applique : \uncertain{l'énoncé} \uncertain{actuel} dans les cas a) et b).

7.3. On peut donner une autre forme précisée de 7.1, en y faisant varier le premier $\ell$, les \uncertain{nombres} premiers \uncertain{(Serre et Tate)} \struck{\ill{} \ill{}} demandent \add{dans} distincts de $p$. \struck{\ill{} \ill{}} \add{$\Gamma$,} \struck{\ill{}} groupe \add{$H^{i}$} \add{pour deux $\ell$} \struck{\ill{} \ill{} distincts} que les \struck{représentations} \add{\uncertain{deux}} associés (quotients de $\pi$) … $= s$, et que les représentations \uncertain{distinctes}, \uncertain{soient} les \og \uncertain{mêmes} \fg{} \uncertain{dans} \uncertain{un} \ill{} $\Gamma$ \uncertain{suivant} \ill{} \uncertain{plus} \struck{\ill{}} … $\}$ de $\mathbf{Q}_{\ell}$, $\mathbf{Q}_{\ell'}$ ; \struck{\ill{}} \struck{$\mathrm{Tr}\, g_{H^{i}(\overline{X}, \mathbf{Q}_{\ell})}$ …} \add{\ill{}} une \uncertain{extension} commune des \uncertain{deux}, on \uncertain{va} \uncertain{dans} \uncertain{les} précisément, \struck{\ill{}} on \uncertain{veut} \struck{\ill{}} \uncertain{plus} \uncertain{généralement} \uncertain{pour} $g \in I$, et qu'il soit indépendant de $\ell \neq p$ :
\[ \text{(7.3.1)} \qquad \mathrm{Tr}\, g_{H^{i}(\overline{X}, \mathbf{Q}_{\ell})} = \mathrm{Tr}\, g_{H^{i}(\overline{X}, \mathbf{Q}_{\ell'})} \]
pour tout $g \in \pi$. [Ceci sera prouvé pour $i = 1$ par \struck{\ill{}} \uncertain{dans} les conditions de Néron.] \struck{De la suite, on montrera \ill{} \ill{} \ill{} \ill{} \ill{} questions \ill{} \ill{} \ill{} Serre et Tate} \ill{} \ill{} \ill{}.

7.4. Lorsque $k$ est fini, (Serre et Tate) \add{\uncertain{(dans $\mathbf{Q}$)}} ont \uncertain{formulé} la conjecture que \uncertain{cette} relation (7.3.1) entraîne \uncertain{encore} $\forall g \in \pi$ dont l'image dans
\[ \mathrm{Gal}(\overline{k}/k) \simeq \widehat{\mathbf{Z}} \]
\struck{\ill{}} \add{\uncertain{tombe}} dans $\mathbf{Z}$. [Cela \uncertain{sera} \uncertain{précisé} … notre \uncertain{séminaire}.]

\marginal{(où $\mathrm{frobenius}$ \struck{\ill{}} les \uncertain{valeurs} \uncertain{propres} \uncertain{des} $g$ \uncertain{sur} $H^{i}(\overline{X}, \mathbf{Q}_{\ell})$ \uncertain{entières} \uncertain{algébriques} \uncertain{de} \uncertain{valeur} \uncertain{absolue} … $q^{w/2}$, où $q = \mathrm{card}\, k$, $0 \leq w \leq 2i$)\note{note marginale au crayon, en bas à gauche, reliée par une accolade au second 7.4 ; lecture très incertaine.}}

\uncertain{En} \uncertain{outre} \uncertain{il} \uncertain{y} \uncertain{a} \uncertain{à} \uncertain{voir} que si $g \in \pi$ va sur $f^{-1}$ \add{(\emph{$f$ le générateur} …)} \note{« $f^{-1}$ » est cerclé ; un mot raturé à sa suite.}

7.4. On peut préciser \uncertain{partiellement} les conjectures \uncertain{en} \uncertain{tenant} compte des \uncertain{représentations} \uncertain{de} \uncertain{tout} \uncertain{le} \uncertain{groupe} \struck{\ill{}} \uncertain{des} groupes algébriques $\Gamma_{\mathbf{Q}_{\ell}} \times \Gamma_{\mathbf{Q}_{\ell}}$, \uncertain{par} \uncertain{exemple} de $\Gamma_{\mathbf{Q}_{\ell}}$,

\page{35}\note{p. 34 de l'auteur.}
en remarquant que ce groupe algébrique sur $\mathbf{Q}_{\ell}$ provient d'un groupe algébrique \struck{\ill{}} $G_{\mathbf{Q}} \times \Gamma_{\mathbf{Q}}$ sur $\mathbf{Q}$. On pourrait demander \add{d'abord} \struck{que} si les représentations algébriques sur les \uncertain{adjoint} $\mathbf{Q}_{\ell}$ \uncertain{en} \uncertain{fait}, par les $H^{i}(\overline{X}, \mathbf{Q}_{\ell})$ proviennent d'une $=$ représentation \uncertain{algébrique} sur $\mathbf{Q}$. \add{ce n'est raisonnable cependant} \struck{(\ill{} \ill{} \ill{} \ill{} \ill{})} \ill{} car nulle, ou une démonstration par voie transcendante \struck{pour $i = 1$ …} (utilisant la cohomologie entière transcendante) devrait être possible. En car.\ résiduelle $> 0$, on peut remarquer que …

\marginal{que ces représentations s'écrivent (de façon \uncertain{pas} \uncertain{unique} à \uncertain{isom.}\ près) \uncertain{sur} $\mathbf{Q}$ ; on \uncertain{peut} \uncertain{demander} qu'elles \uncertain{soient} \uncertain{indépendantes} de $\ell$. [N.B.\ $H$ la représentation \uncertain{induite} \uncertain{ne} \uncertain{s'écrit} \uncertain{pas} \uncertain{sur} $\mathbf{Q}$ \uncertain{pour} $\ell \neq p$, \uncertain{sinon} \uncertain{pour} $i = 1$ et \uncertain{une} \uncertain{variété} \uncertain{abélienne} sur $K$ \uncertain{ayant} \uncertain{« bonne réduction potentielle »}.\note{note de marge écrite en oblique le long du bord gauche, reliée au texte par un trait ; lecture d'ensemble très incertaine.}}

Dans le cas d'un corps résiduel fini, … \og l'\uncertain{homologie} algébrique \fg{} \uncertain{donne} un groupe
\[ \mathbf{Z}_{\mathbf{Q}}(G_{\mathbf{Q}} \times \Gamma_{\mathbf{Q}}) = G \]
sur $\mathbf{Q}$, \uncertain{produit} semi-direct, \add{\uncertain{de générateur} $f_{q} \in \mathbf{Z}_{\mathbf{Q}}(\mathbf{Q})$ …} \uncertain{de} $G_{\mathbf{Q}}$ \struck{\ill{}} par \uncertain{la} multiplication par $q$, et \uncertain{par} $\Gamma$ \uncertain{comme} il \uncertain{opère} sur $(\Gamma/[\Gamma,\Gamma])(\ell)$ \uncertain{il} \uncertain{opère} \uncertain{ainsi} \uncertain{par} \uncertain{multiplication} par $\Gamma$]. On \uncertain{peut} \uncertain{alors} demander que les représentations sur les \uncertain{adjoint} $\mathbf{Q}_{\ell}$ par les $H^{i}(\overline{X}, \mathbf{Q}_{\ell})$ \uncertain{soient} \struck{définies} \add{\uncertain{déduites}} ($\ell \neq p$) \uncertain{définies} par \struck{… $\mathbf{Q}$, et \uncertain{déduites}} d'une $=$ \uncertain{représentation} : celle-ci \uncertain{déduite} … représentation de $G$ sur $\mathbf{Q}$.

\page{36}\note{p. 35 de l'auteur.}
\section{8. Théorème de monodromie : démonstration arithmétique}
\note{titre de l'auteur, souligné.}

\textbf{Th.\ 8.1.} \emph{Supposons que $V$ soit limite inductive \add{filtrante} de sous-anneaux locaux $V_{i}$,} \struck{\emph{l'idéal}} \add{\emph{l'idéal maximal $\mathfrak{m}_{i} = \mathfrak{m} \cap V_{i}$}} \emph{contenant l'uniformisant $t$, tels que les deux conditions suivantes soient satisfaites :}

\emph{a) \struck{le des} $V_{i}$ \uncertain{est} \add{\uncertain{régulier}}} \struck{\emph{et le diviseur} $\mathrm{div}(t)$ \emph{est} … \emph{dans} $\mathrm{Spec}\, V_{i}$ \emph{est à croisements normaux}} \add{\emph{et $V_{i}/tV_{i}$ sont \uncertain{réguliers}}}\struck{\emph{, i.e.\ $t$ fait partie d'un système régulier de param.\ de $V_{i}$}}.

\emph{b) Le corps \add{(résiduel)} $k_{i} = k(V_{i})$ est tel qu'aucune extension finie de $k_{i}$ ne contienne toutes les racines $\ell^{\nu}$-ièmes de l'unité, pour $\ell \neq$ car.\ $k$.}

\emph{Alors le th.\ de monodromie \struck{\ill{}} sous la forme 7.1 est vrai pour \struck{\ill{}} \add{tout} \add{préschéma} $X$ \struck{\ill{}} sur $K$,} \struck{\emph{tel que}} \add{\emph{lorsque l'opération de $I$ sur} \emph{$H^{i}(X_{\overline{K}}, \mathbf{Z}_{\ell})$ $\otimes_{\mathbf{Z}_{\ell}} \mathbf{Q}_{\ell}$ se fait à travers $I_{t}$,}} \struck{$H^{i}(X_{\overline{K}},$ $\mathbf{Z}/\ell\mathbf{Z})$ \emph{soit fini}} \struck{\emph{dès}} \emph{\uncertain{dès} que les $V_{\alpha}$ sont des anneaux de val.\ discrète.}

\textbf{Corollaire 8.2.} \emph{Le théorème de monodromie} \add{\emph{sous la forme 7.1}} \emph{est vrai} \struck{\emph{lorsque}} \struck{\emph{$H^{i}(X_{\overline{K}},$}} \struck{\emph{$\mathbf{Z}/\ell\mathbf{Z})$ est fini}} \emph{pour $X_{K}$ propre sur $K$ dans \struck{\ill{}} chacun des cas suivants :}

\emph{a) $V$ est \add{le hensélisé de, le hensélisé strict de} l'anneau local d'un point d'un schéma $S$ de type fini sur un \add{un localisé de} corps $k$, ou sur \struck{$\mathrm{Spec}$} $\mathbf{Z}$.}

\emph{b) $V$ est de car.\ nulle, (i.e.\ $K$ est de caractéristique} \struck{\emph{résiduelle}} \emph{nulle).}

\textbf{Démonstration 8.2.} \struck{Cas a)} \struck{Supposons $V$ le hensélisé de l'anneau local d'un … $S$ est de type fini sur $\mathrm{Spec}\,\mathbf{Z}$,}

\emph{Cas} a). \struck{Si …} Supposons d'abord que $V$ soit l'anneau \uncertain{de} \uncertain{v.\ d.}\ \struck{\ill{}} \add{qui} \uncertain{est} \uncertain{un} \uncertain{localisé} \struck{\ill{}} l'anneau local d'un schéma $S$ de type fini sur \struck{$\mathrm{Spec}$} \add{un localisé de} $\mathbf{Z}$, ou sur un corps $k_{0}$. Dans le premier cas, \add{(où $S$ est de type fini sur le corps premier)} le corps \struck{\ill{}} résiduel $k(V)$ est de type fini sur le corps premier, donc $V$ satisfait la condition

\page{37}\note{p. 36 de l'auteur ; le bas de la page est vide.}
Dans le cas d'un corps de base $k_{0}$ quelconque, $k_{0}$ \struck{\ill{}} est limite inductive de sous-corps $k_{\alpha}$ de t.f.\ sur le corps premier, et pour $\alpha$ assez grand $S_{\alpha}$ provient d'un $S_{\alpha}$ sur $k_{\alpha}$. Soit $x_{\alpha}$ l'image de $x \in S$ dans $S_{\alpha}$. Comme $\mathcal{O}_{S,x}$ est \uncertain{plat} sur $\mathcal{O}_{S_{\alpha}, x_{\alpha}}$, on a
\[ \dim \mathcal{O}_{S_{\alpha}, x_{\alpha}} \leq \dim \mathcal{O}_{S, x} = 1 , \]
et $\mathcal{O}_{S_{\alpha}, x_{\alpha}}$ est régulier, donc $\mathcal{O}_{S_{\alpha}, x_{\alpha}}$ est un corps \struck{ou \uncertain{spécial} de} ou un anneau de val.\ discrète.

D'ailleurs pour $\alpha$ assez grand, c'est un anneau de valuation discrète, \struck{et $t \in \mathfrak{m}_{\alpha}$ est} (sinon $\mathcal{O}_{S,x}$ serait un corps !), \struck{et pour $\mathfrak{m}_{\alpha}$ est engendré par $t$} et $t \in \mathfrak{m}_{\alpha}$. \struck{Comme $\mathfrak{m}/\mathfrak{m}^{2} = \varinjlim \mathfrak{m}_{\alpha}/\mathfrak{m}_{\alpha}^{2}$, on} \struck{voit que} $t \notin \mathfrak{m}_{\alpha}^{2}$ (\uncertain{puisque} \uncertain{non} \uncertain{dans} $\mathfrak{m}^{2}$), donc on est sous les conditions de a), b) de 8.1.

\page{38}\note{p. 37 de l'auteur. Le haut de la page (jusqu'à « $S_{\alpha}$ ») est barré de grandes croix ; il est transcrit tel qu'il se lit.}
\struck{Dans le cas d'un corps \ill{} \ill{} \ill{} \ill{} \ill{}, $V$ \ill{} écrit} \struck{8.1. Dans le \ill{}, $V$ \ill{} dans un $V \otimes_{k_{0}} k'_{0}$, où $k'_{0}$ est une extension radicielle convenable de $k_{0}$, $k_{0} = \varinjlim k_{\alpha}$, où $k_{\alpha}$ parcourt les sous-corps de $k_{0}$ qui sont des extensions de type fini du corps premier. Pour $\alpha$ assez grand, $S$ provient d'un $S_{\alpha}$ \ill{} \ill{} et le diviseur \ill{} de $S$ provient d'un diviseur \add{réduit} $D_{\alpha}$ sur $X_{\alpha}$. Alors $V$ est la limite des $\mathcal{O}_{S_{\alpha}, x_{\alpha}}$, $x_{\alpha}$ le point générique de $D_{\alpha}$ \ill{} \ill{} considéré dans $S_{\alpha}$. Or le \ill{} $\mathcal{O}_{S_{\alpha}, x_{\alpha}}$ est un anneau de valuation discrète, \ill{} \ill{} de type fini (\ill{} \ill{}).}

\marginal{(en quelques \uncertain{endroits} \uncertain{de} \uncertain{groupes} \uncertain{d'inertie}), \ill{} \uncertain{remarquer} \uncertain{que} \uncertain{si} $k(V)$ \uncertain{est} \uncertain{de} \uncertain{type} \uncertain{fini} sur $k$ … \uncertain{séparable} … $S$ \uncertain{lisse} sur $k_{0}$, et … \uncertain{Écrivons}\note{note de marge écrite en oblique le long du bord gauche, reliée par une accolade au début du paragraphe barré ; lecture très incertaine.}}

\add{\uncertain{Comp.}\ $V \otimes_{k_{0}} k'_{0}$ !}

Dans les deux cas, l'anneau \struck{\ill{}} \add{(par un ou plusieurs d'indices locaux)} est \uncertain{limite} \uncertain{filtrante} \add{de} \struck{\ill{}} \uncertain{sous-anneaux} de valuation discrète, $V_{\alpha}$ à corps résiduel \struck{\ill{} \ill{} \ill{} de type fini. Donc \ill{}} \ill{} \ill{}\supplied{.} … Le hensélisé $V^{h}$ \add{\uncertain{sh}} (resp.\ $(V \otimes_{k_{0}} k'_{0})^{h} = V^{h} \otimes_{k_{0}} k'_{0}$) est limite inductive des $V_{\alpha}^{h}$, et on peut lui appliquer 8.1. De … pour le hensélisé strict.

b) Cas $V$ de car.\ nulle, \add{i.e.\ algébrique sur $\mathbf{Q}$.} ($V$ est limite inductive de \uncertain{ses} \uncertain{sous-$\mathbf{Q}$-algèbres} \struck{$\mathbf{Q}_{i}$} locales $V_{\alpha}$, … de type fini sur $\mathbf{Q}$, \struck{\ill{}} $V_{\alpha}$ d'idéaux maximaux $\mathfrak{m} \cap V_{\alpha} = \mathfrak{m}_{\alpha}$. On peut supposer $t \in V_{\alpha}$. Par Hironaka, on peut trouver un \struck{\ill{}} schéma propre, \struck{\ill{}} intègre, \emph{régulier} $S'_{\alpha}$ sur $S_{\alpha} = \mathrm{Spec}\, V_{\alpha}$, avec $S'_{\alpha} \to S_{\alpha}$ birationnel, et \struck{\ill{}} \add{$\mathrm{div}(t_{S'_{\alpha}})$} l'image inverse \struck{\ill{}} de $\mathrm{div}(t_{S_{\alpha}})$ un diviseur à croisements normaux. \struck{Prenant} Comme $S'_{\alpha} \to S_{\alpha}$ est propre, il existe un point $x'_{\alpha}$ de $S'_{\alpha}$ tel que $V$ \uncertain{domine} l'anneau local $V'_{\alpha} = \mathcal{O}_{S'_{\alpha}, x'_{\alpha}}$ soit dominé par $V$. \struck{Donc} Or c'est un anneau local régulier, et le diviseur de $t$ y est à croisements normaux. \struck{Donc} Les $V'_{\alpha}$ ayant cette propriété supplémentaire

\marginal{\struck{\ill{}} $t = \prod x_{i}^{n_{i}}$, \uncertain{dans} … $x_{i} \in \mathfrak{m}_{\alpha}$ … \uncertain{les} $n_{i} \geq 1$ … \uncertain{supposons} … $n_{1} = 1$\note{note en bas à gauche, en partie encadrée ; lecture très incertaine.}}

\page{39}\note{p. 38 de l'auteur.}
($V_{\alpha}$ régulier, \struck{$\mathrm{div}(t\ \text{dans}\ V_{\alpha})$ \ill{}} $V_{\alpha}/tV_{\alpha}$ régulier$^{*}$) sont cofinaux \uncertain{parmi} les $V_{\alpha}$. De plus, tout $V_{\alpha}$ a un corps résiduel de type fini. On est donc sous les conditions de 8.1.

\textbf{Remarque 8.3.} \struck{Cette dém.} La démonstration du cas b) \struck{\ill{}} serait valable sans restriction sur $V$, une fois qu'on disposerait de la résolution des singularités \add{sous la forme} de Hironaka \struck{\ill{}} (avec diviseur à croisements normaux …) pour les schémas de type fini sur $\mathbf{Z}$ ; \uncertain{il} \uncertain{suffirait} \uncertain{même}, indiquons que \struck{\ill{}} l'hyp.\ de 8.1 sur $V$ est \emph{très} satisfaite. Elle serait en particulier encore \ill{} \ill{} si $V$ est de car.\ $p > 0$, si on dispose de la résolution des singularités : pour les schémas de t.f.\ sur $\mathbf{F}_{p}$. Notons qu'on \uncertain{aurait} \uncertain{seulement} besoin \uncertain{ici} \struck{\ill{}} \struck{\ill{}} que de \og l'\emph{uniformisation locale} \fg.

8.4. \emph{Démonstration de 8.1.} On \struck{\ill{}} \add{peut supposer les $V_{\alpha}$ henséliens} … \uncertain{quitte} à les \uncertain{remplacer} par leurs hensélisés. On a
\[ K = \varinjlim \underbrace{(V_{\alpha})_{t}}_{K_{\alpha}} , \]
donc pour $\alpha$ assez grand, $X_{K}$ provient d'un \struck{\ill{}} \struck{\ill{}} schéma \struck{\ill{}} \struck{vérif.}\ $X_{\alpha}$ sur \struck{$K_{\alpha}$} \add{\uncertain{un}} $(V_{\alpha})_{t}$, i.e.\ \struck{\ill{}} sur l'ouvert $U_{\alpha} = S_{\alpha t}$ de $S_{\alpha} = \mathrm{Spec}\, V_{\alpha}$, \struck{\ill{}} soit $\eta_{\alpha}$ le pt générique de $S_{\alpha}$ i.e.\ de $S_{\alpha t}$, c'est l'image de $\eta_{\alpha}$. \struck{Comme $f_{\alpha} : X_{\alpha} \to U_{\alpha}$ … propre …} \struck{Choisissons …} Comme $f_{\alpha}$ est propre, on sait \add{(\uncertain{SGA} 5 VI)} que
\[ E_{U_{\alpha}} = R^{i} f_{\alpha *}(\mathbf{Z}_{\ell}) \text{ sur } U_{\alpha} \]
est constructible, donc \struck{\ill{}} \uncertain{sur} un voisinage ouvert $U'_{\alpha}$ de $\eta_{\alpha}$ sur lequel $E_{\alpha}$ soit localement constant tordu. \uncertain{Augmentant} au besoin $\alpha$, \add{(… comme $g^{-1}(U_{\alpha}) = g^{-1}(N_{\alpha}) = \{\eta\}$),}

\marginal{$S \ni \eta_{\alpha}$ ; $g : S \to S_{\alpha} \supset U_{\alpha}$\note{petit schéma de l'auteur dans la marge inférieure gauche ; disposition reproduite approximativement.}}

\page{40}\note{p. 39 de l'auteur. Le bas de la page (à partir de « Nous allons utiliser ») est barré de deux longs traits obliques ; l'argument continue au-delà du lot.}
on peut supposer $V_{\alpha} = U_{\alpha}$ (SGA IV 8), donc $E_{\alpha}$ est un faisceau constant tordu, sur $U_{\alpha}$, \add{($E_{\alpha}$ \uncertain{provient} d'une représentation \uncertain{$\ell$-adique} de $\pi_{1}(U_{\alpha}, \overline{\eta})$ dans un …)} \uncertain{et} \uncertain{par} le th.\ de changement de base \uncertain{lisse} \uncertain{propre} (SGA 4 XII), \marginal{donnant $E$ de t.f.\ sur $\mathbf{Z}_{\ell}$}
\[ R f_{K *}\bigl((\mathbf{Z}_{\ell})_{X_{K}}\bigr) = E_{K} \]
est \struck{l'image} isomorphe à l'image inverse de $E_{U_{\alpha}}$ par \struck{$\ldots$} $\eta \to U_{\alpha}$, donc la représentation de $\pi$ \struck{\ill{}} définie par $H^{i}(X_{\overline{K}}, \mathbf{Z}_{\ell})$ est celle déduite de $\pi_{\alpha}$ par l'homomorphisme canonique
\[ \pi \longrightarrow \pi_{\alpha} \qquad \text{\struck{$\ldots \to \pi_{\alpha}$}} \]
\struck{Notons que $\pi_{\alpha}$ est une extension de $\pi_{\alpha,0} = \mathrm{Gal}(\overline{k}_{\alpha}/k_{\alpha})$} Par le \uncertain{même}, \uncertain{remplaçant} l'\uncertain{exposant} $\ell$ au lieu de l'indice $\alpha$, et \uncertain{abrégeant}, on \uncertain{a}
\[ \pi \longrightarrow \pi' . \]
Notons que $\pi'$ est une extension de
\[ \pi'_{0} = \mathrm{Gal}(\overline{k}'/k') \simeq \pi_{1}(S'_{\alpha}) \]
par un ss-groupe \struck{\ill{}} $I'$, et que l'hom.\ \uncertain{induit} un hom.\ :
\[ I \longrightarrow I' \subset \pi' . \]
Dans le cas où $V_{\alpha}$ est de val.\ discrète, \struck{\ill{}} la conclusion résulte de n° 6.

\marginal{Ceci suffit déjà à traiter le cas 8.2 a), qui ne sera \uncertain{utile} que pour donner l'\uncertain{idée} de la dém.\ \og générale \fg{} de l'exposé suivant (\uncertain{comme} \uncertain{ici} \uncertain{à} \uncertain{titre} \uncertain{d'exemple}) … \uncertain{à} \uncertain{traiter} \uncertain{le} \uncertain{cas} \uncertain{des} \uncertain{variétés} \uncertain{abéliennes} … \uncertain{de} \uncertain{choisir} $V$ \uncertain{comme} \uncertain{limite} \uncertain{de} $V_{\alpha}$ … \uncertain{dans} \uncertain{quelques} \uncertain{cas} !)\note{note de marge écrite en oblique le long du bord gauche, reliée au texte par une accolade ; lecture d'ensemble très incertaine.}}

\struck{Nous allons utiliser dans le théorème de \uncertain{structure} … des \uncertain{variantes} du « \uncertain{lemme} d'\uncertain{Abhyankar} » :}

\struck{\emph{Lemme 8.5.} Soit $A$ un anneau local noethérien \uncertain{hensélien}, $k$ \uncertain{son} corps résiduel $k$, … l'exp.\ caractéristique de $k$, $S = \mathrm{Spec}\, A$, $D = \sum_{i \in I} D_{i}$ un diviseur à croisements normaux … d'un syst.\ de paramètres $(x_{i})_{i \in I}$, $x = \prod x_{i}^{n_{i}}$, $U = S_{x}$, $\xi$ un pt géom.\ de $U$. Alors $\xi$ définit une \uncertain{clôture} \uncertain{séparable} $\overline{k}$ de $k$, et …}

\end{document}
